% ---- 01-abstract.md ----
\begin{abstract}

People often explore knowledge by following curiosity---moving from one idea to the next---rather than by asking a single question. We present \textbf{Constellations}, an interactive system for \textbf{low-friction exploratory navigation} that constructs \textbf{bipartite (Atomic\(\leftrightarrow\)Composite)} graphs on demand from natural-language queries. Users begin with a single entity and iteratively expand the graph while the system enforces bipartite alternation. To reduce instability from mid-graph ontology drift, Constellations selects a bipartite pair from the first query (currently among Person\(\leftrightarrow\)Event, Ingredient\(\leftrightarrow\)Recipe, Symptom\(\leftrightarrow\)Disease, Author\(\leftrightarrow\)Paper) and \textbf{locks it for the session} (no switching). To support interpretation and trust, Constellations attaches \textbf{evidence-backed edges} (short snippets with source links) so users can inspect the claimed reason behind a connection at any step. The system's primary data sources are \textbf{Wikipedia/Wikimedia APIs} (Wikipedia + Wikidata), academic metadata APIs (e.g., OpenAlex/Crossref), and an \textbf{LLM}; it does \textbf{not} rely on general internet search for evidence, but it does use \textbf{DuckDuckGo image search} as a fallback when Wikimedia does not provide a suitable image. We describe the system design and position it relative to work on two-mode/affiliation networks and event-centric cultural heritage modeling. We argue that evidence-backed bipartite exploration offers a complementary paradigm to search and static knowledge graphs: a lightweight interface for navigating the frontier between what users know and what they might learn next.

\end{abstract}

% ---- 02-introduction.md ----
\section{Introduction}

People often explore knowledge by following curiosity---moving from one idea to the next---rather than by asking a single question. We call the boundary this process reveals the \textbf{knowledge frontier}---the set of adjacent ideas that are reachable from a current concept but not yet part of one's active memory. Interfaces for search, browsing, and encyclopedic reading support directed lookup, but can be less suited to this low-friction, branching process of "what else is nearby?".

Constellations is designed for this mode. Starting from a single user-provided entity, the system constructs a small, local graph and supports repeated \textbf{expand} operations that reveal additional connected entities. The key design decision is to enforce a \textbf{bipartite alternation constraint}: nodes alternate between \textbf{Atomic} entities (individuals or elementary concepts) and \textbf{Composite} entities (events, works, groups, conditions, or other aggregations). This constraint yields a stable interaction model across domains: atomic nodes are expanded into composites, and composites into atomics. In the project's origin domain---history---this corresponds to a simple idea closely aligned with event-centric cultural heritage modeling: \textbf{events as "meetings"} that bring multiple people together. We generalize this to other domains (e.g., films: people\(\leftrightarrow\)movies; cooking: ingredients\(\leftrightarrow\)recipes; medicine: symptoms\(\leftrightarrow\)diseases) without changing the core UI.

Because the system operates in open-world settings and aims for broad domain coverage, Constellations uses a large language model (LLM) to propose candidate neighbors during expansion. However, exploratory systems that rely on generative models face an immediate interpretability challenge: users naturally ask \textbf{why} a connection exists before deciding whether it is worth pursuing. To address this, Constellations attaches \textbf{edge evidence}: short evidence snippets and source links displayed when an edge is selected. Evidence is treated as support for interpretation and verification, not as proof of causal claims.

This paper contributes an interactive exploration system and a design framing for bipartite low-friction exploration. We position Constellations relative to two-mode/affiliation networks and event-centric modeling.

\subsection{Contributions}

\begin{itemize}[leftmargin=*]

\item \textbf{A bipartite exploration interface} that enforces Atomic\(\leftrightarrow\)Composite alternation as a domain-general interaction primitive for click-to-expand exploration.

\item \textbf{A session-level locking mechanism for bipartite pair selection}: the system chooses an Atomic\(\leftrightarrow\)Composite pairing from the first query and locks it for the session to reduce mid-graph "ontology drift."

\item \textbf{Domains + text input in one interface}: users can start from curated domain seeds or type any query, enabling both guided and open-ended exploration.

\item \textbf{Evidence-backed edges} that make edge inspection a first-class interaction via short snippets and source links.

\item \textbf{An LLM-assisted, on-demand graph construction pipeline} that builds local neighborhoods without requiring a precomputed knowledge graph.

\end{itemize}

\subsection{Why films are an especially good domain (informal motivation)}

Many domains contain bipartite structure, but films provide a particularly clean and rewarding instance: people\(\leftrightarrow\)works with high-quality public metadata, stable naming, and strong Wikipedia coverage. This makes it easy to generate meaningful expansions (directors, co-stars, producers; films in a person's oeuvre) and to attach interpretable evidence. In contrast, domains like current politics are both more dynamic and more sensitive; they require stricter evidence policies and careful treatment of ambiguity and recency.

% ---- 03-related-work.md ----
\section{Related Work}

Constellations sits at the intersection of (1) sensemaking-oriented visualization and interactive graph exploration, (2) two-mode / affiliation networks and event-centric modeling, and (3) LLM-mediated interfaces that must communicate uncertainty and evidence. We briefly summarize relevant foundations and highlight how Constellations differs in scope: it emphasizes \textbf{low-friction exploratory navigation} and \textbf{edge-level interpretability} in an open-world setting, rather than statistical inference over a fully observed two-mode matrix.

\subsection{Sensemaking and exploratory visualization}

Classic work on sensemaking characterizes analysis as an iterative process of foraging, externalizing structure, and revising hypotheses. This perspective motivates interfaces that make it cheap to "try a move," inspect the result, and backtrack. In InfoVis, interaction models for exploration and task taxonomies provide a language for describing what users do (e.g., search, browse, compare, derive) and for structuring evaluations beyond accuracy.

In this frame, Constellations targets tasks that are often under-emphasized in search-centric evaluation: recall, discovery, and orientation during iterative branching. The system's bipartite constraint can be viewed as a design mechanism for reducing cognitive load: it narrows the space of possible moves ("expand to the other side") while remaining expressive across domains.

\subsection{Graph exploration interfaces and node-link visualization guidelines}

Node-link diagrams enable intuitive local reasoning ("what connects to what?"), but scale poorly without interaction, filtering, and progressive disclosure. Design guidance for node-link diagrams emphasizes readability, avoidance of "hairballs," and interaction techniques (hover, highlight, selection, focus+context). Constellations embraces progressive disclosure by default: the graph is small and local, and growth is explicitly user-driven through expansion (including bulk expansion of frontier nodes).

\subsection{Two-mode / affiliation networks (bipartite graphs)}

Two-mode (bipartite) networks are a long-standing model for situations where actors participate in events or belong to groups. Canonical examples include actors\(\leftrightarrow\)movies, authors\(\leftrightarrow\)papers, boards\(\leftrightarrow\)companies, and other membership/affiliation structures. A key methodological lesson is that naively projecting a bipartite network into one mode (e.g., actor--actor) can introduce artifacts; bipartite structure deserves analysis and visualization in its own right.

Constellations adopts bipartiteness not primarily as an analytic device, but as an interaction constraint that simplifies exploration: each step alternates between atomic and composite nodes. Joint spatial displays of affiliation networks (e.g., correspondence-analysis-style embeddings) also provide an important precedent for "map-like" knowledge exploration and inform a future direction of large-scale knowledge cartography.

\subsection{Scholarly collaboration and citation graphs}

"Collaboration graphs" in the scholarly domain---co-authorship networks, citation networks, and related bipartite constructions (authors\(\leftrightarrow\)papers; papers\(\leftrightarrow\)venues; papers\(\leftrightarrow\)keywords)---are a classic and widely studied use case for two-mode graphs. These systems often support search and navigation over a curated academic corpus and can enable tasks like literature discovery, influence tracing, and community/field structure exploration.

Constellations is conceptually adjacent to this tradition and now supports an \textbf{Author\(\leftrightarrow\)Paper} mode backed by an academic-paper corpus API (currently OpenAlex). This enables a classic collaboration-graph workflow (authors, papers, venues) while keeping the same low-friction "expand/backtrack" interaction model. As in other domains, Constellations still avoids general web crawling/search and instead uses a small set of structured sources plus an LLM for open-world neighbor proposal where needed.

\subsection{Event-centric modeling in cultural heritage and museums}

Event-centric cultural heritage modeling emphasizes events as the connective tissue that relates people, places, objects, and documents. The "historical events as meetings" framing aligns strongly with Constellations' original person\(\leftrightarrow\)event concept: an event is any construct that brings multiple people into relation. In Constellations' current design, the system selects a bipartite Atomic\(\leftrightarrow\)Composite pairing from the seed term and \textbf{locks it for the session} to prevent mid-graph switching, while still enabling cross-domain use via different seeds.

\subsection{LLMs, structured outputs, and evidence/provenance}

LLM-based interfaces face a core tension: generative flexibility enables broad coverage, but users need interpretable reasons for what the system shows. Work on provenance and evidence in visualization motivates making "why is this here?" a first-class interaction. Constellations operationalizes this at the edge level: edges carry evidence snippets and a source link, and the UI exposes evidence on selection. This supports user judgment in a low-friction exploration loop without claiming causal conclusions.

\subsection{Constellations' positioning}

Constellations can be read as a pragmatic synthesis of (i) interaction design principles for progressive disclosure in node-link exploration, (ii) the representational clarity of two-mode affiliation networks, and (iii) open-world neighbor proposal via an LLM. Relative to classic two-mode network analysis, Constellations does not assume a fully observed matrix or attempt statistical inference; instead, it constructs small neighborhoods on demand and treats the bipartite constraint primarily as an interaction primitive ("expand to the other side") that limits cognitive load. Relative to static knowledge-graph browsers and search engines, Constellations emphasizes rapid branching and backtracking, and it foregrounds interpretability at the granularity where exploration decisions are made: the edge. Finally, relative to free-form LLM browsing experiences, Constellations introduces two stabilizers---bipartite alternation and a session-level lock of the Atomic\(\leftrightarrow\)Composite pair chosen from the seed---to reduce drift and preserve a coherent local structure during exploration.

% ---- 04-system.md ----
\section{System}

\subsection{Overview}

Constellations is an interactive graph explorer that constructs a small, local graph from a user query and supports repeated expansion. The system enforces a bipartite alternation constraint:

\begin{itemize}[leftmargin=*]

\item \textbf{Atomic nodes}: individual entities or elementary concepts (e.g., Person, Ingredient, Symptom, Player).

\item \textbf{Composite nodes}: aggregations, works, events, groups, conditions, or other "meeting-like" constructs that connect multiple atomics (e.g., Movie, Recipe, Disease, Team, Historical Event/Incident).

\end{itemize}

The UI encodes this alternation with distinct visual forms (e.g., circles vs cards) while keeping interactions consistent across domains.

\subsection{Data flow (high level)}

\begin{enumerate}[leftmargin=*]

\item \textbf{User query}: the user can start from a curated domain seed or type an arbitrary query.

\item \textbf{Context retrieval}: fetch lightweight context (e.g., Wikipedia summary) to help disambiguation and mitigate model knowledge gaps.

\item \textbf{Start-pair classification (locked)}: choose a bipartite pair from the first input (currently one of Person\(\leftrightarrow\)Event, Ingredient\(\leftrightarrow\)Recipe, Symptom\(\leftrightarrow\)Disease, Author\(\leftrightarrow\)Paper), then \textbf{lock it for the entire graph} (no switching).

\item \textbf{Expansion}: call the LLM to propose 8--10 neighbors on the opposite side of the bipartite partition, conditioned on the locked pair.

\item \textbf{Evidence attachment}: each proposed neighbor includes an evidence snippet + page title; selecting an edge reveals the supporting citation.

\item \textbf{Caching}: store nodes and edges (with evidence) in a database to reduce repeated calls and support persistence.

\end{enumerate}

\textbf{Sources (important):} Constellations is intentionally \textbf{multi-source}: it uses \textbf{Wikipedia/Wikimedia APIs} (Wikipedia + Wikidata), \textbf{academic corpora/metadata APIs} (currently OpenAlex, with Crossref/DOI metadata as a fallback), and an \textbf{LLM}. The system does \textbf{not} crawl arbitrary websites or run general internet search for evidence; however, it does use \textbf{DuckDuckGo image search} as a fallback when Wikimedia does not provide a suitable image. Evidence snippets are sourced from Wikipedia page text and/or corpus metadata when available; when the system cannot verify a snippet from available sources, it is shown as missing rather than guessed.

\subsection{Bipartite constraint and "events as meetings"}

The original domain (people\(\leftrightarrow\)events) is motivated by an event-centric view: an event is any construct that brings multiple people into relation. This framing generalizes naturally to other domains by choosing a Composite that aggregates multiple Atomics and for which the inverse membership relation is meaningful (e.g., actors in films; ingredients in recipes).

\subsection{Evidence-backed edges}

Each edge carries structured evidence:

\begin{itemize}[leftmargin=*]

\item \textbf{kind}: currently "wikipedia" (verbatim sentence from a Wikipedia page), "openalex" / "crossref" (corpus metadata), "ai" (model-provided when not verifiable), or "none"

\item \textbf{snippet}: one sentence

\item \textbf{pageTitle + url}: where the snippet is claimed to come from (typically a Wikipedia page)

\end{itemize}

Selecting an edge displays this evidence in the sidebar. The goal is not to "prove" an edge, but to support interpretability and user judgment during exploration.

\subsection{Figures (screenshots)}

The experience is dynamic (the graph reconfigures as nodes are added; users can drag nodes to reshape local structure). A few static screenshots help convey the interaction model:

In Figure 1, expanding a composite (film) reveals connected atomic entities (people) and supporting evidence.

\begin{figure}[t]

\centering

\includegraphics[width=\columnwidth]{\detokenize{../../godfather-brando.png}}

\caption{Figure 1. The Godfather expansion (screenshot)}

\end{figure}

Figure 2 illustrates the timeline view, where temporal metadata provides a second organizing lens alongside the network layout.

\begin{figure}[t]

\centering

\includegraphics[width=\columnwidth]{\detokenize{../../godfather-timeline.png}}

\caption{Figure 2. Timeline view example (screenshot)}

\end{figure}

Figure 3 shows a path-seeking example: finding a bipartite path between two entities.

\begin{figure}[t]

\centering

\includegraphics[width=\columnwidth]{\detokenize{../../John Von Neumann to Geoffrey Hinton.png}}

\caption{Figure 3. Path-seeking example (screenshot)}

\end{figure}

Figure 4 shows the people browser used to seed exploration.

\begin{figure}[t]

\centering

\includegraphics[width=\columnwidth]{\detokenize{../../people.png}}

\caption{Figure 4. Browse People (screenshot)}

\end{figure}

\subsubsection{Cross-domain examples (Atomic\(\leftrightarrow\)Composite generalization)}

Figures 5--7 show the same interaction model applied across domains.

\begin{figure}[t]

\centering

\includegraphics[width=\columnwidth]{\detokenize{../../beef.png}}

\caption{Figure 5. Culinary: Beef (screenshot)}

\end{figure}

\begin{figure}[t]

\centering

\includegraphics[width=\columnwidth]{\detokenize{../../LeBron James.png}}

\caption{Figure 6. Sports: LeBron James (screenshot)}

\end{figure}

\begin{figure}[t]

\centering

\includegraphics[width=\columnwidth]{\detokenize{../../sore throat.png}}

\caption{Figure 7. Medicine: Sore Throat (screenshot)}

\end{figure}

For motion, the repo includes a short demo video: \texttt{/demo.mp4}.

\subsubsection{Demo video (in-browser)}

The embedded video can render as a black frame in some contexts, so we include a representative still frame instead:

\begin{figure}[t]

\centering

\includegraphics[width=\columnwidth]{\detokenize{../../demo-frame.png}}

\caption{Figure 8. Demo frame (click through to watch the video)}

\end{figure}

\href{../../demo.mp4}{Watch the demo video}

\subsection{Interaction design: low-friction branching}

Constellations is optimized for low commitment per step:

\begin{itemize}[leftmargin=*]

\item expanding a node is a single click,

\item backing up (choosing a different node) is immediate,

\item bulk operations allow quick frontier growth (e.g., expand all frontier nodes across the whole graph),

\item when a node's first expansion yields very few neighbors, the UI can automatically request "expand more" once to reduce repeated clicks.

\end{itemize}

This interaction model encourages "try and see" exploration, where users do not need to decide a path "once and for all."

\subsection{Caching and persistence}

To keep exploration responsive and reduce repeated API calls, Constellations caches:

\begin{itemize}[leftmargin=*]

\item \textbf{nodes} (title, type, summaries, images, and classification metadata), and

\item \textbf{edges} (role/label and evidence metadata).

\end{itemize}

Caching supports repeated browsing, saving/loading graphs, and revisiting a previously explored frontier with evidence intact.

\subsection{Implementation notes}

The system is deployed as two primary interactive experiences:

\begin{itemize}[leftmargin=*]

\item \textbf{Standalone Web Application}: A full-canvas exploratory research interface. It supports deep-dive sessions, saved graph management, and a comprehensive "People Browser" for seeding new investigations.

\item \textbf{Chrome Extension (Wikipedia Sidebar)}: A context-aware implementation that lives alongside a user's active browsing. Specifically, when reading a Wikipedia article, the extension constructs a relational graph of the article's core entities. This provides something like a \textbf{physical location for the article in information space}, allowing the user to see the topological "neighborhood" of the topic they are currently reading.

\end{itemize}

While the core ideas are model- and stack-agnostic, a few design choices matter for reproducibility and user experience:

\begin{itemize}[leftmargin=*]

\item \textbf{Graph rendering and interaction}: a force-directed node-link layout supports direct manipulation (dragging nodes), hover highlighting, and edge selection for evidence inspection. Nodes are visually differentiated by partition (Atomic vs Composite) to reinforce alternation.

\item \textbf{Bounded neighborhood expansion}: each expansion requests a small number of neighbors (typically 8--10) to avoid hairball growth and to keep the choice set cognitively manageable.

\item \textbf{LLM output guardrails}: expansions are constrained to named entities on the opposite partition and require per-edge evidence fields (a short snippet plus a source page title/URL). This does not guarantee correctness, but it makes edges inspectable.

\item \textbf{Context-first disambiguation}: before classification/expansion, the system fetches lightweight context (e.g., a Wikipedia summary) to reduce ambiguity and mitigate model knowledge gaps.

\item \textbf{Session-level pair locking}: after the first query, the chosen Atomic\(\leftrightarrow\)Composite pair is locked for the remainder of the graph exploration to reduce mid-graph switching and preserve stable semantics for expansion prompts.

\item \textbf{Caching and persistence}: nodes and edges (including evidence) are cached to reduce repeated LLM calls and to enable saving/loading graphs with evidence intact.

\end{itemize}

% ---- 06-discussion-future-work.md ----
\section{Discussion \& Future Work}

\subsection{Discussion: what Constellations is (and is not)}

Constellations is designed for \textbf{exploration}, not inference. While bipartite network analysis offers statistical models for affiliation data, Constellations operates in an open-world setting with locally constructed neighborhoods and focuses on user experience: recall, curiosity, and sensemaking.

\subsection{LLM use case: connection discovery, not text generation}

An important distinction about the role of AI in Constellations: \textbf{we are not using the LLM to write text for the user}. That is typically the primary concern people have about AI systems---generated content that may be misleading, biased, or fabricated. Instead, Constellations uses the LLM to \textbf{find connections} between entities.

This is conceptually very similar to how a search engine works: in response to a query, the search engine returns a series of links to documents, and some of them may not be relevant. Users are already accustomed to this pattern and have developed effective strategies for handling it---they simply skip over the irrelevant results and focus on the useful ones. The cognitive load and trust calibration required are familiar and manageable.

In our application, irrelevant connections may occasionally appear in the graph. This should not be viewed as a catastrophic failure, but rather as an expected characteristic of exploratory retrieval systems. Users can quickly assess connection relevance through the attached evidence snippets and either dismiss spurious edges or follow promising ones. The interaction model naturally supports this filtering behavior without imposing undue burden.

\subsection{Limitations}

\begin{itemize}[leftmargin=*]

\item \textbf{Open-world ambiguity and name collisions}: many entities share names (e.g., people vs works with the same title). Lightweight context helps, but disambiguation remains imperfect and can propagate errors into expansion.

\item \textbf{Model knowledge gaps and recency}: LLMs can be outdated, and Wikipedia coverage varies; the system can misclassify entities or miss important neighbors without stronger retrieval and verification.

\item \textbf{Evidence is supportive, not definitive}: a single snippet and link is often enough for user judgment, but it does not guarantee the edge is correct. Automated snippet verification and multi-source provenance are not yet implemented.

\item \textbf{Domain variance}: some domains naturally form clean affiliation structures (e.g., actors\(\leftrightarrow\)films), while others exhibit weaker "membership" semantics or higher noise. Session-level pair locking reduces drift but can also constrain legitimate cross-type exploration.

\item \textbf{Scalability and layout stability}: force-directed layouts are effective for small local neighborhoods but can become unstable or visually dense under repeated bulk expansion, motivating multi-resolution views and clustering.

\end{itemize}

\subsection{Future work directions}

\subsubsection{Better expansion ranking (bipartite-aware)}

Use bipartite-inspired heuristics for ranking and diversity:

\begin{itemize}[leftmargin=*]

\item down-weight generic hubs,

\item promote diversity across roles/decades/types,

\item rank by evidence quality and specificity.

\end{itemize}

\subsubsection{Stronger evidence and provenance}

Move from "one snippet" to richer provenance:

\begin{itemize}[leftmargin=*]

\item multiple evidence items per edge,

\item automated verification that the snippet exists in the claimed source,

\item user feedback loops (confirm/reject edges).

\end{itemize}

\subsubsection{Pair learning and "soft locks"}

The current system chooses among a small set of bipartite pairs and locks the choice for the session. A natural extension is to treat pair selection as a probabilistic belief state: maintain a small set of candidate pairings with confidence, allow controlled switching only when confidence crosses a threshold, and preserve previously drawn edges by freezing node partitions once placed.

\subsubsection{Global "map-like" exploration (joint embedding)}

Inspired by joint displays of affiliation networks, a long-term direction is to place a very large two-mode graph into a shared 2D space ("knowledge cartography"), enabling map-like exploration at multiple zoom levels. A practical approach would require multi-resolution embeddings and approximate optimization rather than exact global minimization.

\subsubsection{Domain packs and curatorial modes}

For domains like film and museums:

\begin{itemize}[leftmargin=*]

\item curated constraints on allowed composite types,

\item exhibit-specific "packs" with guided seeds and narratives,

\item touch-screen / installation mode for public spaces.

\end{itemize}

\subsubsection{Scaling limits as a design probe}

Bulk frontier expansion (e.g., expanding all leaf nodes repeatedly) provides a practical way to probe system limits: at what point do layout stability, latency, and evidence quality degrade? Characterizing this "interesting limit" could inform adaptive strategies such as multi-resolution views, clustering, and progressive summarization.

% ---- 07-references.md ----
\section*{References (Working List)}

This is a lightweight, link-forward reference list for the "printable snapshot." For the evolving bibliography, see \texttt{../bibliography.html}.

\subsection*{Visualization / sensemaking}

\begin{itemize}[leftmargin=*]

\item Heer, J., \& Shneiderman, B. (2012). \emph{Interactive Dynamics for Visual Analysis}. \url{https://idl.cs.washington.edu/papers/interactive-dynamics/}

\item Brehmer, M., \& Munzner, T. (2013). \emph{A Multi-Level Typology of Abstract Visualization Tasks}. DOI: \url{https://doi.org/10.1109/TVCG.2013.124}

\end{itemize}

\subsection*{Two-mode / affiliation networks (bipartite)}

\begin{itemize}[leftmargin=*]

\item Borgatti, S. P. \emph{Network Analysis of Two-Mode Data} (PDF). \url{https://www.researchgate.net/profile/Stephen-Borgatti/publication/222212892\_Network\_Analysis\_of\_Two\_Mode\_Data/links/5a0df27d45851541b7079ee9/Network-Analysis-of-Two-Mode-Data.pdf}

\item Faust, K. (2005). \emph{Using Correspondence Analysis for Joint Displays of Affiliation Networks} (PDF). \url{https://sites.socsci.uci.edu/\~{}kfaust/faust/research/articles/faust\_correspondence\_analysis\_2005.pdf}

\item CIDOC CRM Tutorial (Stead, 2008) --- includes "Historical events as meetings". \url{https://www.slideshare.net/slideshow/cidoccrmtutorialorigppt/254963731}

\item CIDOC CRM (site/spec). \url{https://cidoc-crm.org/}

\end{itemize}

\subsection*{Knowledge graphs + retrieval}

\begin{itemize}[leftmargin=*]

\item Vrandečić, D., \& Krötzsch, M. (2014). \emph{Wikidata: A Free Collaborative Knowledge Base}. DOI: \url{https://doi.org/10.1145/2629489.2629492}

\item Auer, S., et al. (2007). \emph{DBpedia: A Nucleus for a Web of Open Data}. DOI: \url{https://doi.org/10.1007/978-3-540-76298-0\_52}

\item Lewis, P., et al. (2020). \emph{Retrieval-Augmented Generation for Knowledge-Intensive NLP Tasks}. arXiv: \url{https://arxiv.org/abs/2005.11401}

\end{itemize}

% ---- 08-acknowledgements.md ----
\section*{Acknowledgements / Author Note}

\textbf{Author.} John Dimm, Lean Software Development (sole proprietorship). Degrees from the University of Washington. Previously worked for 9 years in the Advanced Technology Group at Encyclopaedia Britannica.

\textbf{AI assistance.} This draft and the Constellations software were developed with assistance from large language models (LLMs) used as a programming and writing aid (e.g., generating code suggestions, outlining text, and editing prose). The LLM is not an author and does not take responsibility for the content.

\textbf{Media and data sources.} Node summaries and many images are drawn from Wikipedia and Wikimedia Commons; when Wikimedia lacks a suitable image, the system falls back to DuckDuckGo image search. Any errors or omissions in how sources are attributed are the author's responsibility.
